\section{Introduction}
\label{sec:intro}

RNG~\cite{rng} builds a datacenter fabric as a single layer of top-of-rack routers wired as a random graph. Its routing protocol, Spraypoint, sprays each flow to a random neighbor before pointing it at the destination, which spreads every pair's traffic over dozens of paths. That suits throughput, but much datacenter traffic is small and latency-sensitive (RPCs, key-value lookups, control messages), and for it what matters is the tail of the flow completion time (FCT). Spraying sends even a message between two directly connected racks over three or four hops, and each hop is another queue to wait in and another place to be dropped.

We test a \emph{bypass} that lets small flows skip the spray and take a shortest path, which the paper warns concentrates traffic on direct links:

\begin{quote}
\emph{On an RNG fabric, does routing small flows on shortest paths reduce their tail latency? What does it cost the rest of the traffic, and when does the cost outweigh the gain?}
\end{quote}

We answer with an ns-3 simulation, flow-level computations at the paper's scale, and the paper's analytical models extended to the bypass (\secref{sec:models}), testing four hypotheses:
\begin{description}[leftmargin=2.5em,style=nextline]
  \item[H1 (latency)] The bypass cuts small flows' FCT tail, most between neighbors, which cross one queue instead of two to five (\secref{sec:h1}).
  \item[H2 (no free lunch)] Other traffic barely notices at low small-flow volume, but as the small flows' share of bytes or the load grows, direct links become hotspots and past a crossover the bypass makes things worse (\secref{sec:h2-uniform}, \secref{sec:h2-perm}).
  \item[H3 (reach)] Only about $d/(n-1)$ of router pairs are neighbors, so reaching most small flows means covering every pair, which raises the hotspot risk (\secref{sec:h3}).
  \item[H4 (overhead)] A switch pays one extra table lookup, negligible next to queueing; we argue this but do not measure it (\secref{sec:bypass}).
\end{description}
